\documentclass[10pt,a4paper]{amsart}
\usepackage[margin=19mm]{geometry}
\usepackage{amsmath,amssymb,mathtools}
\usepackage[scr=rsfs,cal=euler]{mathalfa}
\usepackage{xfrac}
\usepackage[unicode,hidelinks,bookmarks=false]{hyperref}
\newcommand{\ie}{i.e.\ }
\newcommand{\cf}{cf.\ }
\newcommand{\resp}{respectively\ }
\newcommand{\Subsection}{\subsection}
\providecommand{\nobreakdash}{\nobreak\hbox{-}}
\setlength{\emergencystretch}{2em}
\multlinegap0pt
\usepackage{bm}
\usepackage{bbm}
\usepackage{dsfont}
\usepackage{xspace}
\usepackage{cancel}
\usepackage{mathtools}
\usepackage{amsthm}
\makeatletter
\@ifundefined{theorem}{\newtheorem{theorem}{Theorem}}{}
\makeatother
\newcommand{\vect}[1]{\ensuremath{\boldsymbol{\mathrm{#1}}}}
\title[Intersection of Reinforcement Learning and Bayesian Optimiza]{Intersection of Reinforcement Learning and Bayesian Optimization for Intelligent Control of Industrial Processes: A Safe MPC-based DPG using Multi-Objective BO}
\author{Hossein Nejatbakhsh Esfahani, Javad Mohammadpour Velni}
\date{Source: arXiv:2507.09864v1 (CC BY 4.0)}
\begin{document}
\maketitle

\noindent\emph{Source and licence.} Intersection of Reinforcement Learning and Bayesian Optimization for Intelligent Control of Industrial Processes: A Safe MPC-based DPG using Multi-Objective BO. Version arXiv:2507.09864v1 (CC BY 4.0), distributed under the Creative Commons Attribution 4.0 International licence, as recorded in the arXiv licence metadata for this exact version and shown on the abstract page https://arxiv.org/abs/2507.09864v1. Full rights notice: licenses/arxiv-2507.09864-CC-BY-4.0.txt.\par\smallskip

\noindent\textit{Editorial scope.} Counted unit: the Theorem whose source label is \texttt{theorem:V\_modif} in the article (source0.tex lines 205--211, source label \texttt{theorem:V\_modif}), delivered together with the article's own complete original proof of it (source0.tex lines 212--237, one \texttt{proof} environment of 26 lines). No mathematics is written, repaired or re-derived: statement and proof are the article's own text line for line. Required context retained uncounted: eq:modif\_V\_finite (lines 162--165); eq:modif\_L (lines 167--169). Every definition, display and label of those lines is retained unchanged. Omitted from this package and not redistributed: 1--161, 166--166, 170--204, 238--552. Nothing inside a retained excerpt is omitted or altered: every retained body is delivered exactly as the article's lines are written, so no omission receipt of any kind accompanies this package. Citations: the retained excerpts carry no citation command at all, so no bibliography entry of the article is transcribed and no reference is redistributed. Reference adaptation: none was needed and none was made; every reference inside the delivered unit resolves inside it, so no unresolved reference or citation is produced. Printed slips kept literal: no printed word, symbol, hypothesis or number of the counted statement or of its proof was altered. Layout and class: the article's own class is \texttt{IEEEtran} with packages including amssymb, enumitem, mathtools, amsthm, graphics, epsfig, algorithm2e, amsmath, amsfonts, mathrsfs, cite, dsfont, siunitx, algorithmic; the wrapper is the lane's pinned \texttt{amsart} editorial wrapper at 10pt on A4, which restates the theorem-like environments the retained text needs and adds the article's own macro definitions that the retained text needs (\texttt{\textbackslash vect}), with extra wrapper packages (bm, bbm, dsfont, xspace, cancel, mathtools). The delivered numbering is the unit's own. Assets: no retained span contains a figure environment, a graphics-inclusion command, a PDF-page inclusion, a table or a TikZ picture; no image file is copied into this package, the package declares no asset dependency and no image file is redistributed with it. Licence: version arXiv:2507.09864v1 of this article is distributed under the Creative Commons Attribution 4.0 International licence, as recorded in the arXiv licence metadata for this exact version and shown on the abstract page https://arxiv.org/abs/2507.09864v1. A full rights notice is delivered at licenses/arxiv-2507.09864-CC-BY-4.0.txt. Characters: every retained excerpt is pure ASCII, so the delivered engine, XeLaTeX with its default Latin Modern text font, reports no missing character.
\begin{align}\label{eq:modif_V_finite}
    &\hat V^{\vect\pi,N}\left(\vect x_k\right)=\mathbb{E}\Bigg[\gamma^N T_{\vect\delta}\left(\vect {\hat x}_N\right)+L\left(\hat{\vect {x}}_0,\vect\pi\left(\hat{\vect {x}}_0\right)\right)\\\nonumber
    &\qquad+\sum_{i=1}^{N-1}\gamma^i L_{\vect\delta}\left(\vect {\hat x}_{i},\vect {\hat u}_{i}\right)\Big|\hat{\vect {x}}_0=\vect x_k,\quad \hat{\vect {u}}_i=\vect\pi\left(\hat{\vect {x}}_i\right)\Bigg],
\end{align}

\begin{align}\label{eq:modif_L}
     L_{\vect{\delta}}(\hat{\vect{x}}_i, \vect{\pi}(\hat{\vect{x}}_i)) = V^\star(\hat{\vect{x}}_i) - \gamma \mathbb{E}[V^\star(\hat{\vect{x}}_{i+1})].
\end{align}

\begin{theorem}\label{theorem:V_modif}
    An Optimal Control Problem (OCP) with the associated value function \eqref{eq:modif_V_finite} can yield the optimal value function $V^\star$ associated with the true MDP, such that the following holds:
    \begin{align}
        \hat V^\star\left(\vect x_k\right)=\min_{\vect\pi}\hat V^{\vect\pi,N}\left(\vect x_k\right)=V^\star\left(\vect x_k\right),
    \end{align}
    and the corresponding optimal policy can capture the true optimal policy $\vect\pi^\star$.
\end{theorem}

\begin{proof}
Let us consider the modified stage cost \eqref{eq:modif_L} and choose the terminal cost $T_{\vect{\delta}}(\hat{\vect{x}}_N) = V^\star(\hat{\vect{x}}_N)$. Now, we apply a telescoping sum to the modified value function \(\hat{V}^{\vect{\pi}, N}(\vect{x}_k)\) in \eqref{eq:modif_V_finite}. Using the modified cost functions, we can express the total cost over the horizon as
\begin{align}\label{eq:modifV}
    &\hat V^{\vect\pi,N}\left(\vect x_k\right)=\mathbb{E}\Bigg[\gamma^N V^\star(\hat{\vect{x}}_N)+L\left(\hat{\vect {x}}_0,\vect\pi\left(\hat{\vect {x}}_0\right)\right)\\\nonumber
    &\qquad+\sum_{i=1}^{N-1}\gamma^i\left(V^\star(\hat{\vect{x}}_i) - \gamma \mathbb{E}[V^\star(\hat{\vect{x}}_{i+1})]\right)\Big|\hat{\vect {x}}_0=\vect x_k\Bigg].
\end{align}
Let us define the modified term $V_m$ as
\begin{align}
    V_m=\sum_{i=1}^{N-1}\gamma^i\left(V^\star(\hat{\vect{x}}_i) - \gamma \mathbb{E}[V^\star(\hat{\vect{x}}_{i+1})]\right).
\end{align}
Under Assumption 2, this modified term then becomes a telescoping sum such that
\begin{align}\label{eq:modif_term}
    V_m=\gamma V^\star(\hat{\vect{x}}_1)-\gamma^N\mathbb{E}[V^\star(\hat{\vect{x}}_N)].
\end{align}
Substituting \eqref{eq:modif_term} in \eqref{eq:modifV}, we have that
\begin{align}
    &\hat V^{\vect\pi,N}\left(\vect x_0\right)=L\left(\hat{\vect {x}}_0,\vect\pi\left(\hat{\vect {x}}_0\right)\right)+\gamma\mathbb{E}[V^\star(\hat{\vect{x}}_{1})]\\\nonumber
&\qquad\qquad\qquad=Q^\star(\hat{\vect{x}}_{0},,\vect\pi\left(\hat{\vect {x}}_0\right))=Q^\star({\vect{x}}_{k},\vect\pi\left({\vect {x}}_k\right)).
\end{align}
It implies that
\begin{align}
    \hat V^\star\left(\vect x_k\right)&=\min_{\vect\pi}\hat V^{\vect\pi,N}\left(\vect x_k\right)\\\nonumber
    &=\min_{\vect\pi}Q^\star({\vect{x}}_{k},\vect\pi\left({\vect {x}}_k\right))=V^\star\left(\vect x_k\right),
\end{align}
and ${\vect{\pi}}^\star\left(\vect x_k\right)=\arg\min_{{\vect{\pi}}} \hat V^\star\left(\vect x_k\right)$.
\end{proof}

\end{document}
